% =====================================================================
%  Trabalho Prático 1 - Recuperação de Informação (SCC0282)
%  Modelo Vetorial e BM25 na coleção Cranfield
%
%  IMAGENS NECESSÁRIAS (subir no Overleaf):
%    usp.jpg                                       -> brasão na capa (raiz)
%    figures/fig1b_preprocessamento_compacto.png   -> Figura 1
%    figures/fig3_curva_precisao_revocacao.png     -> Figura 2
%    figures/fig5_ap_por_consulta.png              -> Figura 3
%    figures/fig4_bm25_grade_MAP.png               -> Figura 4
% =====================================================================

\documentclass[10pt,twocolumn,a4paper]{article}

\usepackage[T1]{fontenc}
\usepackage[utf8]{inputenc}
\usepackage[brazil]{babel}
\usepackage[a4paper,top=2cm,bottom=2cm,left=1.8cm,right=1.8cm,columnsep=0.7cm]{geometry}
\usepackage{mathptmx}
\usepackage[scaled=0.9]{helvet}
\usepackage{courier}
\usepackage{amsmath,amssymb}
\usepackage{graphicx}
\usepackage{booktabs}
\usepackage{makecell}
\usepackage{array}
\usepackage{caption}
\usepackage{microtype}
\usepackage[hidelinks]{hyperref}

\captionsetup{font=small,labelfont=bf,justification=centering}
\captionsetup[table]{skip=3pt}
\setlength{\parskip}{0pt}
\setlength{\parindent}{1.2em}

% Marcação de documento relevante nas tabelas de Top-5
\newcommand{\rel}[1]{\textbf{#1}$^{\bullet}$}

% Tipografia dos títulos de seção
\usepackage{titlesec}
\titleformat{\section}{\normalfont\bfseries\large}{\thesection.}{0.5em}{}
\titleformat{\subsection}{\normalfont\bfseries\normalsize}{\thesubsection}{0.5em}{}
\titlespacing*{\section}{0pt}{10pt}{4pt}
\titlespacing*{\subsection}{0pt}{7pt}{3pt}

% =====================================================================
\begin{document}

% ---------------------------------------------------------------------
% Capa (página única, fora do layout de duas colunas)
% ---------------------------------------------------------------------
\begin{titlepage}
\centering
{\large \textbf{Universidade de São Paulo}\\
Instituto de Ciências Matemáticas e de Computação (ICMC)\par}
\vspace{1.5cm}

\includegraphics[width=0.3\textwidth]{usp.jpg}

\rule{\linewidth}{1pt}
\vspace{0.3cm}
{\Huge \textbf{Modelo Vetorial e BM25 na coleção Cranfield: pré-processamento, parametriza\c{c}\~ao e análise de erros}\par}
\vspace{0.3cm}
\rule{\linewidth}{1pt}
\vspace{2cm}

{\large \textbf{Autores:}\par}
\vspace{0.5cm}
\begin{tabular}{lll}
Victor Silva Botelho & N.\textsuperscript{o} USP: 15645421 & \texttt{victorbotelho@usp.br} \\
João Pedro Barbosa Madeira & N.\textsuperscript{o} USP: 13683038 & \texttt{Joaopedro08madeira@gmail.com} \\
\end{tabular}

\vfill

{\large \textbf{Curso:} SCC0282 - Recuperação da Informação (2026/2)\\
\textbf{Professor:} Prof. Marcelo Garcia Manzato\par}
\vspace{0.5cm}
{\large \today\par}
\end{titlepage}

% ---------------------------------------------------------------------
\section{Introdução}

Sistemas de recuperação de informação precisam ordenar uma coleção inteira a
partir de uma consulta curta e ambígua. Os modelos clássicos resolvem isso de
maneiras distintas: o modelo vetorial representa documentos e consultas como
vetores de pesos TF-IDF e ordena pela similaridade do cosseno, enquanto o BM25
parte do arcabouço probabilístico e pondera cada termo com saturação de
frequência e normalização explícita por comprimento. Embora ambos valorizem
termos raros e controlem a influência da frequência, fazem-no de formas
diferentes, e é justamente essa diferença que este trabalho procura tornar
visível.

O objetivo é construir e avaliar um sistema de recuperação textual sobre a
coleção Cranfield, comparando estratégias de pré-processamento, os dois modelos
de ranqueamento e a parametrização do BM25 e, mais importante, explicar por que
os comportamentos observados ocorrem, com base em exemplos concretos. A
principal conclusão é que, nesta coleção, quase toda a discordância entre os
dois modelos é explicada por um único mecanismo: a forma como cada um normaliza
o comprimento dos documentos. Essa hipótese é levantada a partir da inspeção de
consultas individuais e depois confirmada quantitativamente sobre as 225
consultas.

% ---------------------------------------------------------------------
\section{Técnicas Utilizadas}

\subsection{Pré-processamento}

O texto passa por normalização para minúsculas e tokenização por expressão
regular \verb|[a-z]+|, que é determinística e adequada ao texto técnico ASCII do
Cranfield. A escolha tem um efeito colateral relevante: hifens viram
separadores, de modo que \texttt{boundary-layer-control} gera três tokens e
\texttt{photo-thermoelastic} gera \texttt{photo} e \texttt{thermoelast}, fato
que reaparece na Seção~\ref{sec:reform}. A remoção de stopwords usa a lista
\textit{english} do NLTK (198 formas, expandidas para fragmentos alfabéticos) e
o stemming usa o algoritmo de Porter. As stopwords são removidas antes do
stemming, pois a lista é definida sobre formas plenas e o radical de uma
stopword pode coincidir com o de um termo de conteúdo.

\subsection{Modelo Vetorial}

Adotamos o esquema apresentado em aula (ltc.ltc na notação SMART), aplicado
igualmente a documentos e consultas:
\[
  w_{i,j} = \bigl(1 + \log_{10} f_{i,j}\bigr) \times \log_{10}\!\frac{N}{n_i}.
\]
Cada documento é um vetor de dimensão $|V| = 4292$, com uma coordenada por termo
do vocabulário e praticamente todas nulas. O score é a similaridade do cosseno,
\[
  \mathrm{sim}(d,q) = \frac{d \cdot q}{\lVert d \rVert \, \lVert q \rVert},
\]
em que o numerador mede o alinhamento entre consulta e documento e o denominador
remove o efeito da magnitude. Como normalizamos os vetores para norma unitária
já na indexação, o cosseno se reduz a um produto interno: é o mesmo cálculo,
reorganizado para ser feito de uma vez sobre os 1400 documentos. Note que o idf
vale exatamente zero para um termo presente em toda a coleção, que assim é
eliminado do ranking.

\subsection{BM25}

O BM25 foi implementado explicitamente, termo a termo, sem biblioteca de
ranqueamento:
\[
  \mathrm{score}(d,q) = \sum_{t \in q}
  \mathrm{IDF}(t)\,
  \frac{f_{t,d}\,(k_1+1)}
       {f_{t,d} + k_1\!\left(1-b+b\,\dfrac{|d|}{\mathrm{avgdl}}\right)}.
\]
O quociente cresce com a frequência mas satura em $k_1+1$, ao contrário do tf
logarítmico do modelo vetorial; $k_1$ controla quão rápido essa saturação ocorre
(com $k_1 = 0$ o modelo vira binário) e $b$ interpola entre nenhuma normalização
por comprimento ($b = 0$) e normalização integral ($b = 1$). Somamos a
contribuição de um termo tantas vezes quanto ele ocorre na consulta, o que
equivale a $k_3 \to \infty$. Usamos o IDF de Robertson suavizado,
$\log\!\left(1 + \frac{N-n+0{,}5}{n+0{,}5}\right)$, cujo ``$1+$'' impede valores
negativos para termos muito frequentes.

Vale registrar uma identidade algébrica que explica um resultado nosso: aquela
expressão é idêntica a $\log\frac{N+1}{n+0{,}5}$, ao passo que a forma do BIM
vista em aula é $\log\frac{N+0{,}5}{n+0{,}5}$. As duas diferem apenas por uma
constante aditiva $\log\frac{N+1}{N+0{,}5} \approx 0{,}00036$ em todos os termos.
Por isso, ao testarmos as duas variantes, o MAP mudou apenas na sexta casa
decimal (0,304827 contra 0,304821): a escolha entre elas é irrelevante nesta
coleção, e o motivo é algébrico, não empírico.

\subsection{Decisões de implementação}

\textbf{Índice compartilhado.} Ambos os modelos derivam seus pesos da mesma
matriz esparsa de frequências brutas, garantindo que qualquer diferença venha da
função de ranqueamento e não da indexação. \textbf{Texto indexado:} em 1395 dos
1400 documentos o campo \texttt{.W} já começa com o título; concatenar
\texttt{.T} dobraria o tf dos termos do título, então só o fazemos quando o
título ainda não é prefixo do resumo. \textbf{Desempate determinístico:} scores
iguais são ordenados pelo \texttt{doc\_id} crescente, de modo que duas execuções
produzem rankings idênticos.

% ---------------------------------------------------------------------
\section{Avaliação}

\textbf{Coleção.} Cranfield: 1400 documentos de aerodinâmica, 225 consultas em
linguagem natural e 1837 julgamentos de relevância, obtida dos arquivos
originais da University of Glasgow. Cada consulta tem entre 1 e 39 documentos
relevantes (média 7,2), e nenhuma fica sem relevantes.

Dois detalhes da coleção exigem cuidado. Primeiro, os identificadores
\texttt{.I} das consultas não são sequenciais (001, 002, 004, \dots, 365), mas o
arquivo \texttt{cranqrel} numera as consultas de 1 a 225 pela posição no
arquivo; usar o \texttt{.I} como chave desalinha silenciosamente consultas e
julgamentos. Renumeramos por posição, como faz o \texttt{ir\_datasets}. Segundo,
a escala de relevância do Cranfield é invertida: o grau 1 é o melhor (``resposta
completa'') e o 4 é o pior (``interesse mínimo'').

\textbf{Relevância.} Para as métricas binárias seguimos o enunciado: relevante
$\Leftrightarrow$ grau $\geq 1$; graus $-1$ e documentos não julgados contam como
não relevantes. Para o NDCG mantivemos os graus positivos como relevância
graduada, mas corrigindo a inversão ($1\!\to\!4$, $2\!\to\!3$, $3\!\to\!2$,
$4\!\to\!1$): usar o grau bruto daria o maior ganho ao documento menos
relevante. A diferença não é cosmética: o NDCG@10 do BM25 vale 0,355 com o
mapeamento corrigido e apenas 0,304 com o grau bruto.

\textbf{Métricas.} P@10, R@10, F1@10, MAP, MRR e NDCG@10, definidas como em
aula: $\mathrm{AP}(q) = \frac{1}{|R_q|}\sum_k P@k \cdot \mathrm{rel}(k)$, com
$|R_q|$ contando todos os relevantes (relevantes não recuperados penalizam a
métrica); $\mathrm{DCG}@k = \sum_i \frac{2^{\mathrm{rel}_i}-1}{\log_2(i+1)}$; e
curva precisão-revocação interpolada em 11 pontos por
$P(r) = \max_{r' \geq r} P(r')$. Todas foram implementadas por nós e validadas
contra os valores calculados à mão no exercício do slide 41 da Aula 05.

\textbf{Configuração experimental.} O ranking é calculado sobre os 1400
documentos e os cortes @10 aplicados só na avaliação. Os qrels são usados
exclusivamente para avaliar: nenhuma etapa de indexação, ponderação, escolha de
parâmetros ou reformulação os consulta. A configuração principal usa o
pré-processamento vencedor da Seção~\ref{sec:prep} e $k_1 = 1{,}2$,
$b = 0{,}75$.

\textbf{Verificação da implementação.} Uma versão ingênua em Python puro,
escrita com laços diretamente a partir das fórmulas, reproduz os scores da
versão vetorizada em SciPy com tolerância $10^{-10}$, para os dois modelos e
três combinações de $(k_1, b)$. O modelo vetorial também foi comparado ao
\texttt{TfidfVectorizer} do scikit-learn (correlação de Spearman $> 0{,}95$; não
se espera igualdade exata porque o scikit-learn usa
$\mathrm{idf} = \ln(N/\mathrm{df}) + 1$). São 16 testes, em
\texttt{tests/test\_ir\_tp1.py}.

% ---------------------------------------------------------------------
\section{Resultados Obtidos e Análise}

\subsection{Pré-processamento}
\label{sec:prep}

O ganho é monotônico nas quatro configurações e vale para os dois modelos
(Tabela~\ref{tab:prep} e Figura~\ref{fig:prep}): o MAP do BM25 sobe de 0,269
(linha de base) para 0,283 com stopwords, 0,295 com stemming e 0,305 com ambos,
ganho total de 13,2\%. Os dois efeitos são complementares, e não redundantes,
porque atacam problemas diferentes.

\begin{table}[htbp]
\centering
\scriptsize
\setlength{\tabcolsep}{3.5pt}
\renewcommand{\cellalign}{bc}
\begin{tabular}{@{}lrrrrr@{}}
\toprule
\textbf{Configuração} & \textbf{$|V|$} & \makecell{\textbf{$|d|$}\\\textbf{méd.}} &
\makecell{\textbf{MAP}\\\textbf{vet.}} & \makecell{\textbf{MAP}\\\textbf{BM25}} &
\makecell{\textbf{P@10}\\\textbf{BM25}}\\
\midrule
sem stopw./sem stem. & 7044 & 159,1 & 0,243 & 0,269 & 0,216\\
stopwords            & 6925 &  92,1 & 0,250 & 0,283 & 0,226\\
stemming             & 4403 & 159,1 & 0,262 & 0,295 & 0,226\\
\textbf{stopwords + stem.} & \textbf{4292} & \textbf{92,1} &
\textbf{0,276} & \textbf{0,305} & \textbf{0,234}\\
\bottomrule
\end{tabular}
\caption{Efeito do pré-processamento. $|V|$ é o tamanho do vocabulário e
$|d|$ méd. o comprimento médio dos documentos em tokens.}
\label{tab:prep}
\end{table}

A Tabela~\ref{tab:prep} deixa o mecanismo visível. A remoção de stopwords quase não altera o
vocabulário (7044 $\to$ 6925 termos, pois a lista tem só 198 formas) mas elimina
42\% de todos os tokens (comprimento médio 159,1 $\to$ 92,1): o efeito é sobre o
denominador das normalizações, já que a norma do cosseno e o
$|d|/\mathrm{avgdl}$ do BM25 deixam de ser dominados por palavras sem poder
discriminativo. Já o stemming não remove token algum, mas funde formas
flexionadas e reduz o vocabulário em 38\% (7044 $\to$ 4403): o efeito é sobre a
cobertura do casamento léxico, permitindo que uma consulta com
\texttt{solutions} case com um documento que diz \texttt{solution}. Como um ataca
a normalização e o outro a cobertura, somam-se. Vale notar que o stemming
sozinho (0,295) supera as stopwords sozinhas (0,283), o que é coerente com uma
coleção de textos científicos, onde a variação morfológica é intensa. Todos os
resultados seguintes usam a configuração vencedora (stop + stem).

\begin{figure}[htbp]
\centering
\includegraphics[width=0.92\linewidth]{figures/fig1b_preprocessamento_compacto.png}
\caption{MAP e P@10 nas quatro configurações de pré-processamento.}
\label{fig:prep}
\end{figure}

\subsection{Comparação entre os modelos}

\begin{table}[htbp]
\centering
\scriptsize
\setlength{\tabcolsep}{3.2pt}
\begin{tabular}{@{}lrrrrrr@{}}
\toprule
\textbf{Modelo} & \textbf{P@10} & \textbf{R@10} & \textbf{F1@10} &
\textbf{MAP} & \textbf{MRR} & \textbf{NDCG@10}\\
\midrule
Vetorial & 0,217 & 0,378 & 0,249 & 0,276 & 0,487 & 0,326\\
\textbf{BM25} & \textbf{0,234} & \textbf{0,398} & \textbf{0,266} &
\textbf{0,305} & \textbf{0,539} & \textbf{0,355}\\
$\Delta$ relativo & +8,0\% & +5,0\% & +7,0\% & +10,3\% & +10,7\% & +8,7\%\\
\bottomrule
\end{tabular}
\caption{Métricas agregadas sobre as 225 consultas.}
\label{tab:models}
\end{table}

A Tabela~\ref{tab:models} mostra que o BM25 é superior em todas as métricas, com
vantagem maior no MRR (+10,7\%) e no
MAP (+10,3\%), ou seja, ganha sobretudo no topo do ranking. A curva
precisão-revocação (Figura~\ref{fig:pr}) mostra que a vantagem se mantém em
todos os 11 níveis de revocação, e não apenas em média.

A média, porém, esconde o essencial: o BM25 vence em 137 consultas, perde em 81
e empata em 7. A dispersão da Figura~\ref{fig:ap} confirma que há um número
expressivo de consultas abaixo da diagonal. A pergunta interessante não é qual
modelo tem a maior média, mas o que separa as consultas dos dois grupos.

\begin{figure}[htbp]
\centering
\includegraphics[width=0.94\linewidth]{figures/fig3_curva_precisao_revocacao.png}
\caption{Curva precisão $\times$ revocação interpolada em 11 pontos.}
\label{fig:pr}
\end{figure}

\subsection{Análise por consulta}

\begin{table}[htbp]
\centering
\scriptsize
\setlength{\tabcolsep}{2.2pt}
\begin{tabular}{@{}p{1.15cm}ll ccccc r@{}}
\toprule
\textbf{Caso} & \textbf{Q} & \textbf{Mod} &
\textbf{1\textsuperscript{o}} & \textbf{2\textsuperscript{o}} &
\textbf{3\textsuperscript{o}} & \textbf{4\textsuperscript{o}} &
\textbf{5\textsuperscript{o}} & \textbf{AP}\\
\midrule
BM25 $>$ vet. & q223 & vet. & 1008 & 864 & 1399 & \rel{400} & 1359 & 0,19\\
              &      & BM25 & 1399 & \rel{400} & \rel{1398} & 419 & \rel{1387} & 0,46\\
\addlinespace[1pt]
BM25 $>$ vet. & q59  & vet. & 382 & 383 & 842 & \rel{786} & 1154 & 0,11\\
              &      & BM25 & \rel{785} & 381 & 1154 & 149 & \rel{786} & 0,37\\
\addlinespace[1pt]
vet. $>$ BM25 & q113 & vet. & \rel{265} & \rel{748} & 594 & 1272 & 716 & 0,51\\
              &      & BM25 & 704 & 815 & \rel{748} & 638 & 708 & 0,15\\
\addlinespace[1pt]
vet. $>$ BM25 & q190 & vet. & \rel{390} & \rel{391} & 1008 & \rel{15} & 627 & 0,72\\
              &      & BM25 & \rel{390} & 627 & 856 & 1339 & \rel{15} & 0,38\\
\addlinespace[1pt]
ambos ruins   & q124 & vet. & 941 & 898 & 938 & 1070 & 516 & 0,00\\
              &      & BM25 & 941 & 516 & 898 & 442 & 934 & 0,00\\
\addlinespace[1pt]
ambos ruins   & q13  & vet. & 496 & 903 & 520 & 643 & 313 & 0,01\\
              &      & BM25 & 496 & 903 & 520 & 313 & 440 & 0,01\\
\bottomrule
\end{tabular}
\caption{Top-5 de cada modelo nas seis consultas selecionadas. Identificadores
de documento; \textbf{negrito}$^{\bullet}$ = relevante (grau $\geq 1$).}
\label{tab:queries}
\end{table}

A seleção é automática: entre as consultas com ao menos 4 relevantes, as duas de
maior e menor $\Delta$AP e as duas de menor AP médio. Lendo os Top-5 da
Tabela~\ref{tab:queries} lado a lado, aparece um padrão nítido, e é um padrão de
comprimento de documento.

\textbf{(i) BM25 claramente superior: q223 e q59.} Em q223 (\textit{shear
buckling of unstiffened rectangular plates}), o modelo vetorial coloca no topo os
documentos 1008 (51 tokens) e 864 (29 tokens), ambos irrelevantes e muito
curtos, e só encontra um relevante na 4\textsuperscript{a} posição. O BM25
recupera na 3\textsuperscript{a} e 5\textsuperscript{a} posições os documentos
1398 (115 tokens) e 1387 (120 tokens), ambos relevantes e longos. A causa é a
normalização: o cosseno divide pela norma completa do vetor, o que penaliza
integralmente um documento longo, mesmo quando o comprimento vem de tratar o
assunto com profundidade; o BM25 com $b = 0{,}75$ normaliza apenas parcialmente.
Em q59 ocorre o mesmo: o vetorial promove o documento 382 (21 tokens) e o BM25
encontra o relevante 785 já na 1\textsuperscript{a} posição.

\textbf{(ii) Modelo vetorial superior: q113 e q190.} O mecanismo é o mesmo, com
o sinal trocado. Em q113 o BM25 traz para o topo os documentos 704 (185 tokens),
815 (142) e 638 (138), todos longos e todos irrelevantes, enquanto o vetorial
acerta as duas primeiras posições com os documentos 265 (49 tokens) e 748 (83).
Aqui os relevantes são curtos e específicos, e a normalização integral do cosseno
é exatamente o que se quer. Em q190 o vetorial coloca três relevantes no Top-4 e
o BM25 intercala os documentos 856 (137) e 1339 (130), não relevantes. Em outras
palavras: a normalização parcial do BM25 é uma aposta, que paga bem quando os
relevantes são longos e custa caro quando são curtos.

\textbf{(iii) Ambos insatisfatórios: q124 e q13.} Estas revelam limitações que
nenhum ajuste de ranqueamento resolve. A q13 (\textit{what is the basic mechanism
of the transonic aileron buzz}) tem, na 1\textsuperscript{a} posição de ambos os
modelos, o documento 496, ``\textit{a theory of transonic aileron buzz,
neglecting viscous effects}''. O casamento léxico é quase perfeito, e mesmo assim
o julgamento é $-1$, isto é, não relevante. A q124 é uma consulta composta, com
duas necessidades de informação distintas (escoamento em canais delgados e
estabilidade de cascas cônicas); o documento 941, cujo título praticamente repete
a primeira metade da consulta, também é julgado $-1$. Ou seja, o erro não está no
ranqueamento, mas no descompasso entre sobreposição léxica e relevância julgada:
o avaliador julgou contra a necessidade real, que os modelos não enxergam.

\subsection{Verificando a hipótese do comprimento}
\label{sec:hip}

A leitura acima é uma hipótese construída sobre seis consultas. Testamo-la sobre
as 225. Se a normalização for mesmo o mecanismo dominante, então (a) o Top-10 do
BM25 deve ser em média mais longo que o do vetorial e (b) o BM25 deve levar
vantagem justamente nas consultas cujos relevantes são longos.

Ambas se confirmam. O comprimento médio da coleção é 92,1 tokens e o dos
documentos relevantes é 92,5, praticamente idêntico, ou seja, ser longo não torna
um documento relevante. Ainda assim, o Top-10 do modelo vetorial tem em média
79,4 tokens (bem abaixo da média da coleção) e o do BM25, 98,4 (acima). O
vetorial tem, portanto, \textbf{um viés sistemático para documentos curtos que
não é justificado pelos dados}. E a correlação de Spearman entre $\Delta$AP
(BM25 $-$ vetorial) e o comprimento médio dos relevantes de cada consulta é
$\rho = 0{,}432$ ($p \approx 1\times10^{-11}$, $n = 225$). A hipótese levantada a
partir de seis exemplos resiste, portanto, a um teste sobre a coleção inteira.

\begin{figure}[htbp]
\centering
\includegraphics[width=0.88\linewidth]{figures/fig5_ap_por_consulta.png}
\caption{AP por consulta. Pontos acima da diagonal favorecem o BM25.}
\label{fig:ap}
\end{figure}

\subsection{Variação dos parâmetros do BM25}

\begin{table}[htbp]
\centering
\scriptsize
\setlength{\tabcolsep}{3pt}
\renewcommand{\cellalign}{bc}
\begin{tabular}{@{}lrrrrr@{}}
\toprule
\textbf{$k_1 \backslash b$} & \makecell{\textbf{0}\\\textbf{(MAP)}} &
\makecell{\textbf{0,75}\\\textbf{(MAP)}} & \makecell{\textbf{1}\\\textbf{(MAP)}} &
\makecell{\textbf{0,75}\\\textbf{(P@10)}} & \makecell{\textbf{0,75}\\\textbf{(NDCG)}}\\
\midrule
0,5 & 0,268 & 0,291 & 0,294 & 0,221 & 0,338\\
1,2 & 0,277 & 0,305 & 0,306 & 0,234 & 0,355\\
2,0 & 0,281 & \textbf{0,312} & 0,308 & 0,241 & 0,363\\
\bottomrule
\end{tabular}
\caption{Grade $k_1 \times b$. As três últimas colunas fixam $b = 0{,}75$.}
\label{tab:grid}
\end{table}

\begin{figure}[htbp]
\centering
\includegraphics[width=0.88\linewidth]{figures/fig4_bm25_grade_MAP.png}
\caption{MAP do BM25 na grade de parâmetros. A cor codifica a magnitude; o valor
está impresso em cada célula.}
\label{fig:grid}
\end{figure}

O melhor ponto da grade (Tabela~\ref{tab:grid} e Figura~\ref{fig:grid}) é
$k_1 = 2{,}0$, $b = 0{,}75$, com MAP 0,312 e P@10 0,241. O resultado mais claro é que \textbf{$b$ importa muito mais que $k_1$}:
fixado $k_1$, passar de $b = 0$ para $b = 0{,}75$ melhora o MAP entre 8,6\% e
10,8\%, enquanto variar $k_1$ de 0,5 a 2,0 com $b$ fixo rende no máximo 7\%.
Coerente com a Seção~\ref{sec:hip}, a normalização por comprimento é o fator
dominante nesta coleção. Também há interação entre os parâmetros: para $k_1$
pequeno (0,5 e 1,2) o melhor $b$ é 1,0, mas para $k_1 = 2{,}0$ o melhor é 0,75. A
leitura é que $k_1$ alto deixa o score mais sensível à frequência bruta, e aí uma
normalização total ($b = 1$) passa a penalizar demais documentos longos que
acumulam ocorrências legítimas; um pouco de folga ($b = 0{,}75$) compensa.

\textbf{Efeito de $b$ em uma consulta concreta.} A q67 (``\textit{can series
expansions be found for the boundary layer on a flat plate in a\dots}'', 14
relevantes) é o caso mais didático.

\begin{table}[htbp]
\centering
\scriptsize
\setlength{\tabcolsep}{4pt}
\begin{tabular}{@{}lccccc r@{}}
\toprule
\textbf{$b$} & \textbf{1\textsuperscript{o}} & \textbf{2\textsuperscript{o}} &
\textbf{3\textsuperscript{o}} & \textbf{4\textsuperscript{o}} &
\textbf{5\textsuperscript{o}} & \textbf{AP}\\
\midrule
0,0  & 9 & 72 & 1107 & 155 & \rel{2} & 0,14\\
0,75 & \rel{393} & \rel{180} & \rel{4} & \rel{3} & 306 & 0,49\\
1,0  & \rel{3} & \rel{393} & \rel{180} & \rel{4} & \rel{664} & 0,55\\
\bottomrule
\end{tabular}
\caption{Top-5 da consulta q67 para três valores de $b$ ($k_1 = 1{,}2$).}
\label{tab:beffect}
\end{table}

A Tabela~\ref{tab:beffect} traz os três rankings. Com $b = 0$ não há
normalização alguma e o Top-5 é tomado por documentos de 177,
164, 138, 130 e 116 tokens, dos quais apenas o último é relevante: sem
normalização, o score cresce com o tamanho, pois documentos longos simplesmente
acumulam mais ocorrências dos termos da consulta. Com $b = 0{,}75$ o ranking se
inverte por completo, e os cinco primeiros passam a ter 34, 36, 47, 17 e 91
tokens, sendo quatro deles relevantes. Com $b = 1$ os cinco são relevantes e a AP
vai de 0,14 para 0,55. Os documentos relevantes desta consulta são artigos curtos
e diretamente sobre o tema (``\textit{the boundary layer in simple shear flow
past a flat plate}'', 17 tokens); sem normalização eles são esmagados por textos
longos que mencionam os mesmos termos de passagem.

\subsection{Reformulação de consultas}
\label{sec:reform}

Cinco consultas foram reescritas manualmente, cada uma com uma operação diferente
e com a hipótese registrada no código antes de rodar os modelos. As duas versões
de cada consulta foram executadas com os dois modelos, e a
Tabela~\ref{tab:reform} resume as mudanças no Top-10: as colunas P@10 medem
diretamente quantos relevantes entraram ou saíram dessa faixa, e o NDCG@10
acrescenta o efeito de reordenação dentro dela. O Top-10 completo das dez
execuções está em \texttt{results/05\_query\_reformulation\_detail.json}.

\begin{table*}[htbp]
\centering
\small
\setlength{\tabcolsep}{9pt}
\begin{tabular}{@{}llccccc@{}}
\toprule
\textbf{Cons.} & \textbf{Operação} & \textbf{AP vetorial} & \textbf{P@10 vetorial} &
\textbf{AP BM25} & \textbf{P@10 BM25} & \textbf{NDCG@10 BM25}\\
\midrule
q219 & sinônimo              & 0,04 $\to$ 0,12 & 0,00 $\to$ 0,10 & 0,04 $\to$ 0,09 & 0,00 $\to$ 0,10 & 0,00 $\to$ 0,09\\
q69  & + específica          & 0,02 $\to$ 0,05 & 0,00 $\to$ 0,00 & 0,04 $\to$ 0,09 & 0,00 $\to$ 0,10 & 0,00 $\to$ 0,17\\
q109 & + específica (curta)  & 0,03 $\to$ 0,04 & 0,00 $\to$ 0,00 & 0,04 $\to$ 0,06 & 0,00 $\to$ 0,10 & 0,00 $\to$ 0,03\\
q171 & $-$ termos            & 0,31 $\to$ 0,32 & 0,30 $\to$ 0,30 & 0,64 $\to$ 0,56 & 0,30 $\to$ 0,30 & 0,66 $\to$ 0,61\\
q15  & vocabulário           & 0,62 $\to$ 0,58 & 0,20 $\to$ 0,20 & 1,00 $\to$ 1,00 & 0,20 $\to$ 0,20 & 0,74 $\to$ 1,00\\
\bottomrule
\end{tabular}
\caption{Efeito das reformulações no Top-10 (original $\to$ modificada).}
\label{tab:reform}
\end{table*}

\begin{table}[htbp]
\centering
\scriptsize
\setlength{\tabcolsep}{4pt}
\renewcommand{\cellalign}{bc}
\begin{tabular}{@{}lcccc@{}}
\toprule
\textbf{Pos.} & \makecell{\textbf{Vetorial}\\\textbf{original}} &
\makecell{\textbf{Vetorial}\\\textbf{modificada}} &
\makecell{\textbf{BM25}\\\textbf{original}} &
\makecell{\textbf{BM25}\\\textbf{modificada}}\\
\midrule
1  & 1221 & 670        & 1221 & 315\\
2  & 1222 & \rel{666}  & 208  & 1221\\
3  & 208  & 537        & 1222 & 993\\
4  & 299  & 238        & 992  & \rel{666}\\
5  & 993  & 668        & 993  & 1082\\
6  & 664  & 535        & 299  & 371\\
7  & 530  & 1221       & 371  & 139\\
8  & 228  & 139        & 530  & 208\\
9  & 971  & 993        & 7    & 9\\
10 & 197  & 1222       & 1309 & 710\\
\bottomrule
\end{tabular}
\caption{Top-10 da consulta q219 antes e depois da troca de \texttt{small} por
\texttt{low}. \textbf{Negrito}$^{\bullet}$ = relevante (grau $\geq 1$).}
\label{tab:top10reform}
\end{table}

\textbf{q219, sinônimo.} O caso mais limpo. A consulta original diz
``\textit{\dots when the reynolds number is small}'', mas a literatura consagrou
a expressão ``\textit{low Reynolds number}''. Trocando uma única palavra,
\texttt{small} por \texttt{low}, a AP do vetorial vai de 0,045 para 0,119 e a do
BM25 de 0,037 para 0,093, e a P@10, que era zero nos dois modelos, passa a 0,10.
A Tabela~\ref{tab:top10reform} mostra o efeito no Top-10: na versão original,
nenhuma das dez posições traz um relevante nos dois modelos, e as listas são
quase as mesmas (1221, 1222, 208, 299, 993 e 530 aparecem em ambas); na versão
modificada o documento 666 (grau 2) entra em 2\textsuperscript{o} no vetorial e
em 4\textsuperscript{o} no BM25, com sete das dez posições renovadas no vetorial
e seis no BM25. Nada
mudou no significado da pergunta; mudou apenas o vocabulário. É a ilustração
mais direta da limitação apontada na Aula 04: modelos puramente léxicos não
enxergam sinonímia.

\textbf{q69 e q109, tornar mais específica.} Ambas confirmam a hipótese, com
ganho modesto. Em q69 os termos originais são todos pouco discriminativos
(\texttt{solution}, \texttt{boundary}, \texttt{layer} e \texttt{equation} ocorrem
em 400 a 470 dos 1400 documentos); acrescentar \texttt{series}, \texttt{laminar}
e \texttt{compressible} leva a AP do BM25 de 0,039 a 0,092 e a R@10 de 0 a 0,20.
Em q109, uma consulta de apenas 4 tokens, explicitar a consequência física
(tensão térmica, flambagem, flutter) leva a P@10 de 0 a 0,10.

\textbf{q171, remover termos: hipótese parcialmente refutada.} Prevíamos que
enxugar a consulta de 17 tokens beneficiaria sobretudo o modelo vetorial,
prejudicado pela diluição da norma. O vetorial de fato melhorou um pouco
(0,306 $\to$ 0,319), mas o BM25 piorou (0,639 $\to$ 0,556). A explicação é que o
BM25 não normaliza o vetor de consulta: termos adicionais só somam evidência, e
mesmo termos de IDF baixo contribuem positivamente quando presentes. Remover
termos, para o BM25, é descartar evidência sem contrapartida, o que revela uma
assimetria entre os modelos que não havíamos previsto.

\textbf{q15, vocabulário técnico e o valor da métrica graduada.} O termo
\texttt{photoelastic} ocorre em um único documento, mas os dois relevantes usam
\texttt{photo-thermoelastic}, que o tokenizador quebra em \texttt{photo} e
\texttt{thermoelast}, e o stemming não aproxima essas formas. Adotando o termo
dos autores, a AP do BM25 não muda (1,00 $\to$ 1,00), porque os dois relevantes
já estavam no topo; mas o NDCG@10 sobe de 0,74 para 1,00, pois a reformulação
reordenou os dois relevantes, trazendo para a 1\textsuperscript{a} posição o de
grau 1 (``resposta completa''). É um bom argumento a favor de reportar métricas
graduadas: a AP é cega a uma melhora que o usuário perceberia.

\subsection{Análise de erros}

\textbf{Falsos positivos.} Escolhemos documentos não relevantes trazidos pelo
BM25 para as duas primeiras posições. (1) Na q98, o documento 638
(``\textit{longitudinal aerodynamic characteristics at low subsonic speeds of a
highly swept wing utilizing nose flaps}'', grau $-1$) aparece em
1\textsuperscript{o} com score 35,1. A decomposição do score mostra de onde ele
vem: \texttt{flap} ($+6{,}74$), \texttt{control} ($+5{,}70$), \texttt{trail}
($+4{,}87$), \texttt{subson} ($+3{,}03$). Todos os termos da consulta estão lá,
mas a consulta pergunta se controles na ponta se comparam a flaps de bordo de
fuga, e o documento trata de flaps de nariz. É a hipótese de independência entre
termos falhando: o modelo soma evidências isoladas e não percebe que a consulta
pede uma relação entre elas. (2) Na q124, o documento 941 tem título quase
idêntico a metade da consulta e recebe $+8{,}40$ só de \texttt{channel}
($n = 17$, IDF altíssimo), mas é julgado $-1$, porque a necessidade real de
informação está na outra metade da consulta. (3) Um terceiro caso, na q139, é de
outra natureza: os documentos 847 e 897 estão em 1\textsuperscript{o} e
2\textsuperscript{o} e parecem genuinamente pertinentes, mas simplesmente não
foram julgados. Como o enunciado manda tratar não julgados como não relevantes,
eles contam como erro, o que é, na verdade, uma limitação de cobertura dos qrels
do Cranfield, não do sistema.

\textbf{Falso negativo.} Na q80 (``\textit{are methods of measuring aerodynamic
derivatives available which could be adopted for use in short running time
facilities}''), o documento 596 tem grau 3 e aparece na posição 1206 em ambos os
modelos. O motivo é categórico: a interseção entre os tokens da consulta e os do
documento é vazia. O documento se chama ``\textit{the properties of crossed
flexure pivots\dots}'' e descreve o aparato usado na medição, enquanto a consulta
descreve o método e a finalidade. Score zero não é um erro de ponderação: é o
limite do casamento léxico. Nenhum ajuste de $k_1$, $b$ ou pré-processamento
recupera esse documento; seria preciso representação semântica (embeddings) ou
expansão de consulta. Os outros dois falsos negativos que examinamos têm a mesma
assinatura: no documento 236 da q87 (grau 4), o único termo casado é
\texttt{flow}, o de menor IDF da consulta ($n = 730$ de 1400).

% ---------------------------------------------------------------------
\section{Considerações finais}

\begin{enumerate}\setlength{\itemsep}{2pt}
\item \textbf{O pré-processamento é barato e vale a pena.} Stopwords e stemming
juntos elevam o MAP do BM25 em 13,2\%, e são complementares porque atacam frentes
distintas: um corrige o denominador das normalizações, o outro amplia a cobertura
do casamento léxico.

\item \textbf{O BM25 vence, mas não por ser ``melhor'' em abstrato.} Ele supera o
vetorial em todas as métricas (MAP +10,3\%) e em todos os níveis de revocação,
mas perde em 81 das 225 consultas. A diferença é governada pela normalização por
comprimento: o cosseno normaliza integralmente e adquire um viés para documentos
curtos (Top-10 médio de 79,4 tokens contra 92,1 da coleção), enquanto o BM25
normaliza parcialmente. Confirmamos isso com $\rho = 0{,}43$ entre $\Delta$AP e o
comprimento dos relevantes.

\item \textbf{$b$ importa mais que $k_1$}, e os dois interagem: o melhor $b$ cai
de 1,0 para 0,75 quando $k_1$ cresce de 1,2 para 2,0.

\item \textbf{O teto dos dois modelos é léxico.} As consultas em que ambos falham
não falham por ranqueamento: ou o documento relevante não compartilha nenhum
termo com a consulta, ou a consulta é composta, ou o julgamento não acompanha a
sobreposição léxica. Trocar ``\textit{small}'' por ``\textit{low}'' em q219 quase
triplicou a AP, um ganho que nenhum ajuste de parâmetro produziria. É exatamente
onde entram expansão de consulta e representações densas.

\item \textbf{A escolha da métrica muda o diagnóstico.} A AP não viu a melhora da
reformulação de q15, que o NDCG@10 registrou (0,74 $\to$ 1,00); e o NDCG varia
17\% conforme se corrija ou não a escala invertida do Cranfield. Decisões de
avaliação merecem tanta atenção quanto decisões de modelagem.
\end{enumerate}

% ---------------------------------------------------------------------
\section{Uso de ferramentas de IA}

Ferramentas de IA generativa (Claude) foram utilizadas como apoio nas seguintes
atividades: (i) \textbf{apoio à programação}, na estruturação dos módulos, na
vetorização com matrizes esparsas e na escrita dos testes automatizados; (ii)
\textbf{depuração}, no diagnóstico de uma instalação corrompida do SciPy causada
pelo limite de 260 caracteres de caminho do Windows e na identificação do
desalinhamento entre os identificadores \texttt{.I} das consultas e a numeração
do \texttt{cranqrel}; (iii) \textbf{explicação de conceitos}, na discussão das
variantes de IDF do BM25 e das convenções de ganho do NDCG; e (iv)
\textbf{correção ortográfica} deste relatório. As fórmulas, as decisões de
implementação, o desenho dos experimentos, as reformulações manuais de consulta e
a interpretação dos resultados foram revisados e são de responsabilidade dos
autores, que compreendem integralmente o código entregue.

% ---------------------------------------------------------------------
\begin{thebibliography}{9}
\footnotesize
\bibitem{manning}
MANNING, C. D.; RAGHAVAN, P.; SCHÜTZE, H.
\textit{Introduction to Information Retrieval}.
Cambridge University Press, 2008.

\bibitem{robertson}
ROBERTSON, S.; ZARAGOZA, H.
\textit{The Probabilistic Relevance Framework: BM25 and Beyond}.
Foundations and Trends in IR, 2009.

\bibitem{salton}
SALTON, G.; WONG, A.; YANG, C. S.
A Vector Space Model for Automatic Indexing.
\textit{Communications of the ACM}, v. 18, n. 11, 1975.

\bibitem{cleverdon}
CLEVERDON, C. W.
The Cranfield tests on index language devices.
\textit{Aslib Proceedings}, 1967.

\bibitem{jarvelin}
JÄRVELIN, K.; KEKÄLÄINEN, J.
Cumulated Gain-based Evaluation of IR Techniques.
\textit{ACM TOIS}, v. 20, n. 4, 2002.
\end{thebibliography}

\end{document}
